\subsection*{2026-06-23 (the classical gravity field equation, parameter-free: the entanglement first law $\delta S=\delta\langle K\rangle$ with TFPT's atoms gives the LINEARISED Einstein equation with $G$ from $\cthree$; the thermodynamic and geometric origins of $\cthree$ coincide via $|\mu_4|=|\Z_2|\chi=4$; the two formerly-open pieces (matter flux J3, entropy-density coefficient) closed/atom-fixed: \texttt{v358})}
\begin{itemize}
\item \textbf{\texttt{v358} (\texttt{GRAV.ENTROPY.EQUILIBRIUM.01}) --- the bridge to the dynamic world, instantiated on the
gravity side.} Carrying out the Jacobson-2015 / Faulkner-et-al.\ ``first law of entanglement $\Rightarrow$ Einstein
equation'' with TFPT's atoms. \mE\ PARAMETER-FREE COEFFICIENT: $\delta S=\delta\langle K\rangle$ gives
$G_{ab}=\cthree^{-1}T_{ab}$ with $\cthree^{-1}=8\pi$ fixed; all four $\cthree$-roles consistent (Unruh
$1/(2\pi){=}4\cthree$, EH $1/(16\pi){=}\cthree/2$, Einstein $8\pi{=}1/\cthree$, Bekenstein $1/4{=}1/|\mu_4|$) ---
no free dimensionless Newton dial. \mE\ THERMO $=$ GEO COINCIDENCE (new): the first law's coefficient $2\pi/\eta$
equals the geometric $|\Z_2|2\pi\chi$ iff $|\mu_4|=|\Z_2|\chi=4$ (\vref{v216}) --- the thermodynamic and geometric
origins of $\cthree$ are ONE. \mC\ J3 SOLVED: the CHM ball modular Hamiltonian (boost via BW, \vref{v323}) gives the
matter flux $\delta\langle K_B\rangle=\tfrac{8\pi^2R^4}{15}\delta\langle T_{00}\rangle$. \mE\ ENTROPY DENSITY
ATOM-FIXED: $1/4{=}1/|\mu_4|$ (\vref{v57}) $+$ central charge $c=\gcar{+}\Nfam=8$ (form \vref{v59}). \mO\ residual:
linear$\to$full non-linear (the original B6) $+$ the $v_{\mathrm{geo}}$ unit. Closes the LINEARISED Einstein
equation parameter-free; cited in \texttt{tfpt\_research\_contracts}; mirrored in Wolfram (\texttt{302/302}).
\end{itemize}

\subsection*{2026-06-23 (the two frontier walls attacked directly: Direction A assembles the continuum chain of \texttt{SEAM.EQUIV.01} theorem-by-theorem and ties its residual to the seam one-sidedness; Direction B was found ALREADY closed by \texttt{v40} (harmonic metric $=$ finite linear algebra, not a PDE): \texttt{v356})}
\begin{itemize}
\item \textbf{\texttt{v356} (\texttt{SEAM.EQUIV.CONTINUUM.03}) --- Direction A, the continuum chain assembled.} After the
reverse-audit/bandwidth work showed no hidden \emph{discrete} lever remains, the only fundamental progress is to
assemble the \emph{continuum} chain. Link-by-link: \textbf{S1} the gapped ($\Delta{=}6\ln(3/2){>}0$, \vref{v302})
quasi-free $16$-Majorana collar (\vref{v155}/\vref{v160}) $\to$ \textbf{S2} invertible (gapped free fermions have no
intrinsic topological order, Kitaev, \vref{v301}) $\to$ \textbf{S3} the \emph{nontrivial} invertible phase
($c_-{=}8{\neq}0$) $\to$ \textbf{S4} a \emph{non-gappable} anomalous chiral edge (anomaly inflow / bulk--edge), so
edge \emph{existence} is a CONSEQUENCE of S2$+$S3, not a separate assumption --- this directly reduces the
\vref{v336}/\vref{v351} ``does the collar have a chiral edge?'' residual $\to$ \textbf{S5} its chiral-CFT scaling limit
(MMST, \vref{v336}; range $\operatorname{rank}\le c\le D$, $8\le8\le16$) $\to$ \textbf{S6} the order-$4$ $\mu_4$ clock
$+$ the AQFT stack pin $(E_8)_1$ (\vref{v351}/\vref{v297}). \mE\ the only non-theorem link is S3, and \mC\ S3 is the
seam being \emph{one-sided} (chiral, $c_-{\neq}0$) --- the \emph{same} one-sidedness that makes $\cthree{=}1/(8\pi)$ a
one-sided Gauss--Bonnet ($|\Z_2|\oint K$, \vref{v58}). \mO\ the one remaining input is ``the abstract collar is
realized as a genuine lattice chiral free-fermion invertible phase'' (the \vref{v297} Flat-Away). NET: \texttt{SEAM.EQUIV.01}
reduced to ONE realization input tied to the constant that \emph{defines} the theory; it ADVANCES, it does not close.
\item \textbf{Direction B --- found ALREADY closed by \vref{v40}.} The U-wall ``transcendental Hitchin solve'' (\vref{v31})
is the worst case for a \emph{generic} Higgs bundle; but at the \emph{polystable} point TFPT selects, Mehta--Seshadri
makes the bundle \emph{unitary}, so the harmonic metric is the constant invariant Hermitian form (\emph{finite linear
algebra}, Higgs $\Phi{=}0$) and $\det R{=}8$ is read off --- not a PDE. Only the absolute amplitude $U_{\mathrm{point}}$
(an anchor) remains. A redundant new B module was built, found to duplicate/contradict \vref{v40}, and \emph{deleted}
(no padding). Honest outcome: B needs no new work.
\end{itemize}

\subsection*{2026-06-22 (bandwidth search of the unmapped $E_8$ region, with a strict forced-identity discriminator: the genuine overlooked structure is COLLECTIVE ($\sum\deg=128=\dim S^+$, $\prod\deg=|W(E_8)|$, exponents $=$ totatives of $30$); the per-degree atom matches are unforced and DECLINED; no new physical hit -- reconciling \texttt{v66} and \texttt{v354}: \texttt{v355})}
\begin{itemize}
\item \textbf{\texttt{v355} (\texttt{E8.UNMAPPED.BANDWIDTH.01}) --- the next step after \vref{v354}, run honestly.} The
obvious follow-up (``search those five degrees for hits we overlooked'') is itself the numerology temptation
\vref{v354} warns about: a rich object like $E_8$ returns a numerical match for almost anything. So the scan uses a
STRICT discriminator -- a candidate counts only if it is a FORCED identity (a theorem about $E_8$), never a
coincidence. \mE\ FORCED, SET-LEVEL (the genuinely overlooked structure, collective not per-number): $\sum\deg=128=
\dim S^+$, the spinor half of $248=120+128$ (theorem: $\sum\deg=\#$pos roots$+$rank$=120+8$) -- the invariant budget
of $E_8$ equals the fermion content the carrier reads; $\sum\exp=120=\#$pos roots; $\prod\deg=|W(E_8)|=696{,}729{,}600$;
exponents $=$ the $\varphi(30){=}8$ totatives of $30$ (\vref{v223}), so the unmapped degrees' exponents are part of the
forced set; $\max\deg=30=h(E_8)$. \mC\ SUGGESTIVE, not a theorem: $12{=}h(E_6)$, $18{=}h(E_7)$, $30{=}h(E_8)$ are all
degrees of $E_8$ (the $2T/2O/2I$ McKay ladder, \vref{v219}), but ``$\deg(E_8)$ contains its subalgebras' Coxeter
numbers'' is not a general theorem -- flagged, not claimed. \mO\ DECLINED (the numerology temptation): the per-degree
atom matches $14{=}\dim G_2$, $20{=}\det L$, $24{=}|W(A_3)|$, $12{=}|R(A_3)|$, $18{=}p_4(a)$ (\vref{v66} already flagged
$12,18$ as fittable) each admit $\ge2$ readings, and the lone extra prime in $|W(E_8)|=2^{14}3^5 5^2\,7$ ``$=$'' the
scalaron exponent is one of many readings of $7$ -- all declined. \mE\ RECONCILES \vref{v66} (the 8 degrees coincide
with the load-bearing integers) and \vref{v354} (only $2,8,30$ are primary readouts): the search finds NO new forced
physical hit, and the disciplined decline IS the anti-numerology result. (This also sharpens \vref{v354}: the five are
``no PRIMARY readout / hull overhead'', not literally ``unmapped'' -- they coincide with structural atoms, but
unforced.) Cited in \texttt{tfpt\_5\_redteam} (with a full audited degree table); mirrored in Wolfram (\texttt{301/301}).
\end{itemize}

\subsection*{2026-06-22 (the REVERSE numerology audit, and what the golden ratio means: of $E_8$'s 8 Casimir invariants only $3$ feed a primary readout and $5$ are hull overhead; the golden ratio is UNMAPPED structure, so it cannot be numerology: \texttt{v354})}
\begin{itemize}
\item \textbf{\texttt{v354} (\texttt{E8.REVERSE.AUDIT.01}) --- the reverse complement of \vref{v305}.} TFPT's forward
discipline (\vref{v305}) is one-directional: every physical READOUT must map onto an $E_8$ structure (anti-fitting).
The sharp adversarial question runs the OTHER way --- how much $E_8$ structure maps onto NOTHING, and is the
recurring golden ratio numerology? \mE\ THE REVERSE AUDIT: of $E_8$'s eight Casimir invariant degrees
$\{2,8,12,14,18,20,24,30\}$ (exponents $+1$), exactly THREE feed a readout --- $2$ (metric), $8$ (rank $\to\cthree
{=}1/(8\pi)$), $30$ (Coxeter $\to\gcar{=}5$ and the order-$30$ cycle) --- and FIVE carry no primary readout
$\{12,14,18,20,24\}$ (they coincide with structural atoms per \vref{v66}, but unforced -- reconciled in \vref{v355});
so $3/8$ are primary readouts, $5/8$ hull overhead. \mE\ ECONOMY, NOT PROMISCUITY: a small fixed set ($\operatorname{rank}8$,
$h{=}30$, $\det K{=}1$, $D_5{\oplus}A_3$, the $16$ half-spinor) generates the ENTIRE readout set, and the higher
Casimirs are never reached into to fit --- a numerological use would be PROMISCUOUS (mine the rich structure), TFPT is
ECONOMICAL (few invariants, many readouts) and leaves most of $E_8$ untouched, the anti-numerology signature being
precisely the large unused remainder. \mE\ THE GOLDEN RATIO IS UNMAPPED $\Rightarrow$ NOT NUMEROLOGY: $\varphi=2\cos
(\pi/5)$ appears only in internal structure (the affine-$E_8$ attractor eigenvalue \vref{v312}, the icosian lattice
\vref{v348}) and in NO physical readout; numerology means a number FITTED to an observation, $\varphi$ is fitted to
nothing --- it is the algebraic signature of the $5$-fold ($h{=}30{=}2\cdot3\cdot5$), the carrier announcing it is
icosahedral. \mO\ the $5/8$ unmapped is the HULL OVERHEAD ($E_8$ is the minimal even-unimodular hull, not the world
group); where a future higher-invariant readout could live, or genuine over-capacity --- stated plainly, not hidden.
Cited in \texttt{tfpt\_5\_redteam}; mirrored in Wolfram (\texttt{300/300}).
\end{itemize}

\subsection*{2026-06-22 (the bird's-eye rethink: TFPT is a unique self-consistent LOOP with ZERO free dials, not a linear axiom system; $\pi$ is a fixed geometric constant, not a dial; the one residual is the continuum closure: \texttt{v353})}
\begin{itemize}
\item \textbf{\texttt{v353} (\texttt{TFPT.SELFLOOP.01}) --- refining \vref{v352}.} Stepping all the way back: TFPT is
not a LINEAR theory (axioms $\to$ theorems) but a CLOSED SELF-CONSISTENT LOOP whose outputs fix its own inputs.
\mE\ it is a loop: $\gcar{=}5{=}$ max prime of the output $h(E_8){=}30$, and the $8$ in $\cthree{=}\operatorname{rank}
E_8$ -- the outputs fix the inputs (\vref{v6}), with a unique fixed point (\vref{v56}); ``which is the axiom?'' is
the wrong question for a loop. \mE\ ZERO free adjustable dimensionless dials: every load-bearing number is
forced/over-determined. \mE\ $\pi$ is a FIXED geometric constant (the boundary $2$-sphere, $\oint K{=}4\pi$), not a
dial. \mC\ this refines \vref{v352}'s ``framework $+$ $\pi$'': $\pi$ is not a free input and the framework is a
SELECTION PRINCIPLE (the unique self-consistent loop), not a free meta-axiom -- the honest characterisation is ``a
unique self-consistent loop with zero free dials''. \mO\ the one residual is the CONTINUUM closure of the loop (the
chiral-edge existence, \vref{v336}/\vref{v351}); structural-realism (loop complete) vs constructivism (continuum
open) is a philosophical, not numerical, fork -- no free numbers either way. A capstone synthesis.
\end{itemize}

\subsection*{2026-06-22 (two consequences: the $c{=}8$ which-net ambiguity resolved by the order-4 clock, and BOTH axioms shown bootstrap-forced so the only irreducibles are the framework $+$ $\pi$: \texttt{v351}, \texttt{v352})}
\begin{itemize}
\item \textbf{\texttt{v351} (\texttt{SEAM.EQUIV.CONTINUUM.02}) --- the which-net ambiguity falls.} The
long-standing ``$c{=}8$ ambiguity'' ($E_8$ det 1 vs $SO(16){=}D_8$ det 4, both $c{=}8$, flagged open in
\vref{v277}/\vref{v344}) is RESOLVED, non-circularly, by the seam's ORDER-4 $\mu_4$ clock: the index-4 simple-current
extension of $D_5\!\oplus\!A_3$ is $E_8$ (\vref{v154}), the index-2 would be $SO(16)$, so the order-4 clock (the
four Gauss--Bonnet marks, \vref{v216}) selects $E_8$ (full $\mu_4$ condensation $16{\to}4{\to}1$) not $SO(16)$
(partial $\Z_2$ $16{\to}4$). The bootstrap agrees ($h(E_8){=}30$ max prime $5{=}\gcar$ vs $h(D_8){=}14$ max prime
$7$). \mE\ so $\det K{=}1$ is a consequence, not an open input. \mO\ the residual of \texttt{SEAM.EQUIV.01} is then
PURELY the chiral-edge / scaling-limit existence (\vref{v336}), not which-net.
\item \textbf{\texttt{v352} (\texttt{TFPT.IRREDUCIBLE.01}) --- the inputs are not axioms.} The deepest answer to
``the inputs are back-determined'': BOTH P1 and P2 are bootstrap-forced fixed points. \mE\ $\gcar{=}5$ is forced
three ways (\vref{v6}); the $8$ in $\cthree$ four ways ($\operatorname{rank}E_8=h(D_5)=\phi(30)=\det R$, \vref{v54});
the $E_8$ hull is forced (unique even unimodular rank-8, \vref{v83}/\vref{v154}). \mE\ so the ONLY irreducibles are
the FRAMEWORK (a discrete reflection-positive boundary with a finite carrier) and the transcendental $\pi$. \mO\ the
framework is the META-AXIOM --- not math-reducible, the foundational stance (like ``spacetime is a manifold''),
minimal, tested by its consequences not proved. TFPT's true input is ONE framework $+$ $\pi$. Wolfram mirror $+2$
(297 $\to$ 299).
\end{itemize}

\subsection*{2026-06-22 (correction: the inputs are bootstrap-forced fixed points, not free axioms -- $g_{\mathrm{car}}{=}5$ IS the icosahedral $5$ of $h(E_8){=}30$, and the golden ratio is emergent from the $\mu_4$-glue, not external: \texttt{v350})}
\begin{itemize}
\item \textbf{\texttt{v350} (\texttt{SEAM.EQUIV.BOOTSTRAP.01}) --- correcting \vref{v349}.} An honest correction
prompted by the right objection that the inputs are not axioms but are back-determined by the theory. \mE\ the
inputs ARE bootstrap-forced (\vref{v6}): $\gcar{=}5$ is over-determined three ways -- rank-fill ($\gcar{+}\Nfam{=}8$),
Coxeter-match ($\gcar{=}$ max prime of $h(E_8){=}30$), Pascal ($2^4{=}C(5,0){+}C(5,1){+}C(5,2)$). \mE\ so
$\gcar{=}5$ IS the icosahedral $5$-fold (the $5$ in $h(E_8){=}30{=}2{\cdot}3{\cdot}5$), not the $D_5$ rank; the atoms
$(2,3,5)$ ARE the primes of the icosahedral Coxeter number $30$. \mE\ the golden ratio is EMERGENT from the glue:
$D_5$ ($h{=}8$)/$A_3$ ($h{=}4$) are non-golden alone, but $D_5\!\oplus\!A_3\to E_8$ ($\mu_4$ index-4, \vref{v154})
gives $h(E_8){=}30$ (golden) -- the seam's own $\mu_4$ produces the $5$-fold, so \vref{v349} checked the un-glued
pieces, the wrong object, and its pentagon-vs-count dichotomy was false. \mO\ the residual relocates: the golden is
bootstrap-forced, the only open thing is the physical continuum realisation (\vref{v336}), NOT the golden ratio.
Honest caveat: self-consistency within the framework, not from-nothing. Wolfram mirror $+1$ (296 $\to$ 297).
\end{itemize}

\subsection*{2026-06-22 (the honest test -- does the raw seam carry $\varphi$? -- answer NO: the bare carrier gives $\sqrt2$, golden is the icosahedral input; the keystone reduces to ``is $g_{\mathrm{car}}{=}5$ a pentagon or a count?'': \texttt{v349})}
\begin{itemize}
\item \textbf{\texttt{v349} (\texttt{SEAM.EQUIV.GOLDEN.01}).} The honest, computed test of the \vref{v348}
question, and a clarifying NEGATIVE: the raw seam does \emph{not} carry the golden ratio. \mE\ the golden ratio
$2\cos(\pi/5)$ is the signature of a $5$-fold rotation ($5\mid h$); the carrier $D_5$ has $h{=}8$ (eigenvalue
$\sqrt2$) and $A_3$ has $h{=}4$ -- no 5-fold; the dynamical data (cusp $\{0,1/3,2/3\}$, recovery $(2/3)^6$, seed
$\varphi_0$) are rational/transcendental. \mE\ $5\mid h$ holds only for $A_4{=}SU(5)$ ($h{=}5$), $H_3$ ($h{=}10$),
$E_8$ ($h{=}30$) -- the output/icosahedral side, so $\varphi$ is the icosahedral INPUT. \mE\ $\varphi$ alone would
not suffice ($2I\neq C_5\times\mu_4$, $20\neq120$). \mO\ the resolution: the keystone reduces to one question --
is $\gcar{=}5$ a PENTAGON (a $5$-fold rotation $\Rightarrow$ golden $\Rightarrow$ icosahedral $2I=E_8$) or just a
COUNT (the rank of $D_5\Rightarrow\sqrt2$)? The ``absurdly simple solution'' is real but is a \emph{reading} of
axiom~P2 (the golden ratio is the pentagon content of $\gcar{=}5$, not derivable from $\gcar{=}5$-as-a-count); there
is no hidden $\varphi$. L2 is foundational: P2-as-pentagon closes it (mode~C), else a rigidity theorem must produce
the 5-fold (mode~B); not closed. Wolfram mirror $+1$ (295 $\to$ 296). A NEGATIVE/clarification; prevents a false
closure.
\end{itemize}

\subsection*{2026-06-22 (Route B carried out: rigidity $+$ the icosian-ring identity reduce the whole keystone to ONE question -- does the raw seam carry the golden ratio? -- the simplest form of L2: \texttt{v348})}
\begin{itemize}
\item \textbf{\texttt{v348} (\texttt{SEAM.EQUIV.RIGID.01}).} Executes the rigidity route (mode B of \vref{v347}) and
lands on a strikingly simple statement. \mE\ rigidity closes the crystallographic part: the order-4 clock has the
unique normal form $z\mapsto iz$ (Nielsen/Kerckhoff, \vref{v180}), Troyanov gives the unique flat pillowcase metric
(\vref{v284}). \mE\ the downstream is a LATTICE identity, not just a graph: by Conway--Sloane the ring of icosians
($120$ unit icosians $=2I$, golden-ratio quaternion norm) IS the rank-$8$ $E_8$ lattice. \mE\ the single bridge is
the golden ratio: the crystallographic restriction allows only orders $\{1,2,3,4,6\}$ (the $5$-fold absent), so
$\varphi=2\cos(\pi/5)$ is exactly what extends the crystallographic $\mu_4$ to the non-crystallographic icosahedral
$2I$ (the quasicrystal). \mO\ so the entire keystone reduces to ONE question: \emph{does the raw seam carry the
golden ratio?} -- if yes, $\mu_4{+}\varphi\Rightarrow2I\Rightarrow E_8$ and rigidity fixes the rest. \mO\ honest
subtlety: TFPT's $\varphi$ is currently $E_8$-side (\vref{v312}/\vref{v313}); the raw-seam data so far are not
manifestly golden, so showing the RAW seam carries $\varphi$ non-circularly is the residual ($=$ L2, sharpest form).
Wolfram mirror $+1$ (294 $\to$ 295). An analysis/reduction; closes nothing.
\end{itemize}

\subsection*{2026-06-22 (can the one open arrow be \emph{fully} solved? the closure-mode classification: the flat$\to$spherical base locus, four modes, no free theorem -- only rigidity or axiom $P_3$: \texttt{v347})}
\begin{itemize}
\item \textbf{\texttt{v347} (\texttt{SEAM.EQUIV.CLOSURE.01}).} A direct, honest answer to ``is there any other way
to fully solve the one open arrow L2?''. \mE\ the precise locus: the seam gives the \emph{flat} pillowcase base
$S^2(2,2,2,2)$ ($\chi_{\mathrm{orb}}=2-4(1-\tfrac12)=0$, v216) while the $2I/E_8$ structure is the Seifert base over
the \emph{spherical} $S^2(2,3,5)$ ($\chi_{\mathrm{orb}}=2-(\tfrac12+\tfrac23+\tfrac45)=\tfrac1{30}$), so L2 bridges
two geometric types -- why it needs the carrier $(2,3,5)$ input. \mE\ as a ``physics-realises-geometry'' statement
L2 has exactly four closure modes: A construct (the analytic wall), B rigidity (a unique-realisation theorem --
the only theorem-route w/o new input, pursued v284/v300/v331), C axiom $P_3$ (then TFPT is complete relative to 3
axioms), D empirical. \mE\ the ``no new state'' discipline (v246) blocks deriving L2 from extra physics, so no route
makes it a free theorem. \mE\ verdict: full closure is either a rigidity theorem (B) or axiom status (C), validated
either way by (D); only B closes it as a theorem. \mO\ NOT closed -- classifies the modes, locates the difficulty.
Wolfram mirror $+1$ (293 $\to$ 294). An analysis/roadmap; closes nothing.
\end{itemize}

\subsection*{2026-06-22 (the geometric bridge made explicit: the keystone is ONE arrow -- raw-seam $\Rightarrow\det K{=}1$ is a six-link chain, five theorems $+$ one open orbifold realisation, the group forced by the $(2,3,5)$ atoms: \texttt{v346})}
\begin{itemize}
\item \textbf{\texttt{v346} (\texttt{SEAM.EQUIV.GEOM.01}).} The honest capstone of the \texttt{SEAM.EQUIV.01}
investigation: it lays out the full path raw-seam $\Rightarrow\det K{=}1$ as six links and shows exactly one is open.
\mE/\mC\ L1 raw seam $\to\mu_4$ $+$ four marks (Gauss--Bonnet, \vref{v216}); \mO\ L2 $\mu_4$ $+$ the $(2,3,5)$ atoms
$\to$ the $2I$ orbifold $\mathbb{C}^2/\Gamma$ (THE ONE OPEN ARROW $=$ the geometric content of \texttt{SEAM.EQUIV.01});
\mE\ L3 $2I\to$ affine $E_8$ McKay (\vref{v219}); \mE\ L4 $2I\to E_8$ du Val (\vref{v232}); \mE\ L5
$\mathbb{C}^2/2I\to$ Poincar\'e link $H_1{=}0$ (\vref{v345}); \mE\ L6 $H_1{=}0\to\det K{=}1$. \mE\ five of six links
are theorems; \mE\ the group is FORCED by the seam's own atoms (the exceptional $SU(2)$ axis signatures
$(2,3,3)/(2,3,4)/(2,3,5)\to 2T/2O/2I$; $(2,3,5)$ unique to $2I$, and $(|\Z_2|,\Nfam,\gcar){=}(2,3,5)$ ARE those
axes); \mE\ everything downstream of L2 is then forced. \mO\ the irreducible open content is L2 alone (``the raw
seam normal bundle is an ADE orbifold point $\mathbb{C}^2/\Gamma$'') $=$ \texttt{SEAM.EQUIV.01}; not closed, but the
bridge is now demonstrably ONE arrow with the group pinned and all downstream proven.
\end{itemize}

\subsection*{2026-06-22 (executing route R3: the $(2,3,5)$-rewrite attractor link computed to be the Poincar\'e homology sphere ($\det K{=}1$), the combinatorial core complete, only the geometric bridge open: \texttt{v345})}
\begin{itemize}
\item \textbf{\texttt{v345} (\texttt{SEAM.DETK.02}).} Carries out the combinatorial core of the hypergraph-homotopy
route (R3, the most promising fresh angle from \vref{v344}) as far as it computably goes. By plumbing topology
(Brieskorn/Hirzebruch) the link of a tree plumbing with the ADE Cartan intersection form has
$H_1=\operatorname{coker}(\mathrm{Cartan})$. \mE\ the Smith normal form of the $E_8$ Cartan matrix is the identity,
so $\operatorname{coker}(E_8)=0$ -- the $E_8$-graph plumbing boundary is an INTEGRAL HOMOLOGY SPHERE (the genus-1
meaning of $\det E_8{=}1$); among ADE only $E_8$ does this ($A_n\!\to\!\Z/(n{+}1)$, $D_n\!\to\!\Z/4$,
$E_6\!\to\!\Z/3$, $E_7\!\to\!\Z/2$). \mE\ $\pi_1=2I=SL(2,5)$ is PERFECT (commutator subgroup $=$ the whole order-120
group, computed directly), so $H_1=0$ and the link is THE Poincar\'e sphere, $|2I|=120=|\mu_4|\,h(E_8)$. So R3's
combinatorial half is \mE-complete: the $E_8$ attractor link is forced to be the Poincar\'e sphere ($\det K{=}1$),
uniquely among ADE. \mO\ what stays open is only the geometric realisation bridge (the raw rewrite seam attractor
\emph{is} this $E_8$ du~Val plumbing) $=$ the \texttt{SEAM.EQUIV.01} content; not closed, but the keystone is now
reduced to that single geometric identification. Wolfram mirror $+1$ (292 $\to$ 293).
\end{itemize}

\subsection*{2026-06-22 (the one keystone bit $\det K{=}1$ mapped: six equivalent faces, the genus-1 obstruction, three access routes, the hypergraph-homotopy route the most promising fresh angle: \texttt{v344})}
\begin{itemize}
\item \textbf{\texttt{v344} (\texttt{SEAM.DETK.01}).} A bird's-eye, full-spectrum map of \texttt{SEAM.EQUIV.01}'s
only open residual, the holomorphy bit $\det K{=}1$. \mE\ it is ONE statement with six equivalent faces
(holomorphy / homology-sphere link $H_1{=}0$ / one $1$-dim irrep / $\det$ Cartan $=1$ / full $\mu_4$ condensation /
$\tau{=}i$ CM), unified by $|\det\mathrm{Cartan}|=|H_1(\text{link})|=\#(\text{1-dim irreps})$ across ADE
($A_n{\to}n{+}1$, $D_n{\to}4$, $E_6{\to}3$, $E_7{\to}2$, $E_8{\to}1$) -- only $E_8$ gives $1$. \mE\ the GENUS-1
OBSTRUCTION explains why it is hard: every easy argument (plane vacuum, gap, free-fermion) is genus-0, necessary
but not sufficient (\vref{v87}); $\det K$ is the torus degeneracy. \mE\ THREE genus-1 access routes: R1
continuum-modular (\vref{v336}), R2 anyon-condensation (\vref{v281}), R3 hypergraph-homotopy. \mC\ the HYPERGRAPH
route is the most promising fresh angle: the $(2,3,5)$ rewrite attractor is the affine-$E_8$ graph (\vref{v312}) and
the $E_8$ du~Val link is the Poincar\'e sphere (\vref{v232}), so $\det K{=}1$ becomes a finite combinatorial
statement (no continuum limit). \mO\ none closes it; the gap is the combinatorial-to-geometric bridge $=$ the
\texttt{SEAM.EQUIV.01} content. A synthesis/roadmap; closes nothing.
\end{itemize}

\subsection*{2026-06-22 (the four black-hole-birth completion routes investigated: none closes all or is 100\% provable alone; all four converge on \texttt{SEAM.EQUIV.01} $+$ the decoupled QG contour: \texttt{v343})}
\begin{itemize}
\item \textbf{\texttt{v343} (\texttt{FOUR.ROUTES.01}).} A direct, honest answer to ``which black-hole-birth route
closes ALL problems and is 100\% provable'': \emph{none individually} --- but the four converge on the same two
residuals as \vref{v335}. \mE\ \textbf{(A)} the finite causal diamond reframes \texttt{QG.AMB} to a finite thermal
trace (de Sitter $S_{\mathrm{dS}}=32\pi^4 e^{2\ainv}$ finite, $\chi=729/665$) but inherits the GHP contour
(\vref{v334}, \mO); does not touch the seam. \mC\ \textbf{(B)} Carlip--Cardy gives a second route to existence$+c$
($c{=}8$ consistent, but the Carlip ambiguity and no holomorphy $\det K{=}1$) --- inherits \texttt{SEAM.EQUIV.01}.
\mE\ \textbf{(C)} the modular/thermal flow (\vref{v239}) is well-defined but its geometric content is
\texttt{QGEO.SYM.01}, a corollary of \texttt{SEAM.EQUIV.01} (\vref{v335}) --- a reduction. \mE/\mO\ \textbf{(D)} the
unique attractor (\vref{v56}) gives parameter-freeness but the cyclic dynamics is \mC\ interpretation. \mE\
\textbf{Verdict:} independent closures $=0$; irreducible residuals after all four $=2$ (unchanged from \vref{v335});
the horizon premise makes both more natural but eliminates neither --- focus $=$ \texttt{SEAM.EQUIV.01}. An analysis
module; closes nothing, fabricates nothing.
\end{itemize}

\subsection*{2026-06-22 (the heat-kernel origin of the $\alpha$ terms: the Calder\'on quadratic is the Gilkey gauge-curvature coefficient; the cubic Maxwell moment stays the \texttt{EM.WARD.01} residual: \texttt{v342})}
\begin{itemize}
\item \textbf{\texttt{v342} (\texttt{EM.WARD.02}) --- the heat-kernel sharpening of \texttt{EM.WARD.01}.} Derives the
ORIGIN and STRUCTURE of the $F_{U(1)}$ determinant-line terms from textbook Seeley--DeWitt / Gilkey coefficients,
advancing the EM-Ward origin from \mC\ (conjectured) to \mE\ for the term structure; no new gate. \mE\ $\cthree=
1/(8\pi)$ \emph{is} the boundary Gauss--Bonnet coefficient ($\oint_{S^2}K=2\pi\chi=4\pi$, $\chi{=}2$, \vref{v216}).
\mE\ the $\alpha^2$ (Calder\'on) term \emph{is} the Gilkey gauge-curvature coefficient ($30\,\Omega^2/360=\tfrac1{12}
\Omega^2$, $\Omega{=}i\alpha F\Rightarrow-(\alpha^2/12)F^2$), so it is forced, not an ansatz. \mE\ the $\cthree$-orders
$\{0,3,6\}$ are an arithmetic ladder and $2\cthree^3=1/(256\pi^3)$ carries $\pi^3$ (three boundary insertions). \mC\
the exact coefficient needs the seam $F$-normalisation. \mO\ the cubic $\alpha^3$ Maxwell moment (a topological/Chern
level) $+$ the full ``this is the seam $U(1)$ $\zeta$-determinant'' identification stay the \texttt{EM.WARD.01}
residual, tied to \texttt{SEAM.EQUIV.01}. $\ainv$ stays \mE\ (\vref{v3}). Wolfram mirror $+1$ (290 $\to$ 291).
\end{itemize}

\subsection*{2026-06-22 ($\alpha$ as the stationarity of a $U(1)$ determinant line: every ingredient an index / heat-kernel coefficient / discriminant form, the open EM-Ward obligation named: \texttt{v341})}
\begin{itemize}
\item \textbf{\texttt{v341} (\texttt{ALPHA.QUILLEN.01}).} Answers the external review's sharpest valid point ---
``derive the \emph{form} of $\alpha$, not more decimals''. The EM closure
$F_{U(1)}=\alpha^3-2\cthree^3\alpha^2-\tfrac45\cthree^6 M\log(1/\varphi_{\mathrm{seam}})=0$ is rewritten as a
\emph{stationarity} of the boundary $U(1)$ determinant line:
$[\alpha^3-2\cthree^3\alpha^2]_{\,\delta\log\det\Delta_{U(1)}}=[8b_1\cthree^6]_{\,\text{counterterm}}\cdot
[\log(1/\varphi_{\mathrm{seam}})]_{\,\text{seam response}}$. \mE\ every ingredient is a NAMED object, not a knob:
$M{=}41$ is a $U(1)$ index ($\tfrac45 M=8b_1$, $b_1{=}41/10$ the SM $U(1)_Y$ beta, EM Ward \vref{v48});
$\cthree{=}1/(8\pi)$ the Gauss--Bonnet boundary coefficient ($8{=}\operatorname{rank}E_8$); $-5/4=-q(D_5)$ a
discriminant form and $\varphi_{\mathrm{base}}{=}1/(6\pi){=}\tfrac43\cthree$ carries the other,
$\cthree/\varphi_{\mathrm{base}}{=}3/4{=}q(A_3)$; $\delta_{\mathrm{top}}{=}\Omega_{\mathrm{adm}}\cthree^4$. The
Quillen split is verified at the $\alpha^{-1}{=}137.0359992$ root. \textbf{No new gate:} $\ainv$ stays closed
\mE; v341 adds no open item --- it \emph{sharpens} the pre-existing \mO/physical-origin obligation
\texttt{EM.WARD.01} (``why \emph{this} function'', already conjectural). \mO\ the one still-open part \emph{is} that
residual: the from-first-principles AQFT proof that $F_{U(1)}$ is the EXACT Quillen functional (the variational
principle derived, not assembled), unchanged in scope. Wolfram mirror $+1$ (289 $\to$ 290).
\end{itemize}

\subsection*{2026-06-22 ($A_s$ reconciled with the archived derivation: predicted \& dimensionless (so it does not define a mass), the tension is the reheating channel: \texttt{v340})}
\begin{itemize}
\item \textbf{\texttt{v340} (\texttt{AS.RECONCILE.01}).} Answers ``why don't we have $A_s$ / the old versions
derived it'' and ``can a ratio define mass itself''. \mE\ $A_s$ \emph{is} predicted and \emph{dimensionless}:
$A_s=N_\star^2(M_{\mathrm{scal}}/\bar M_{\mathrm{Pl}})^2/(24\pi^2)=N_\star^2\cthree^7/(24\pi^2)$ with
$M_{\mathrm{scal}}/\bar M_{\mathrm{Pl}}=\cthree^{7/2}=1.26\times10^{-5}$ a pure ratio, so $A_s$ pins the e-fold
count $N_\star$, NOT $\bar M_{\mathrm{Pl}}$ --- it cannot define an absolute mass (No-Unit, \vref{v153}, stands).
\mC\ the $-11.4\sigma$ at the slow Higgs channel ($N_\star{=}51.4$) is the slow-vs-fast reheating dichotomy;
the $A_s$-preferred $N_\star{=}56.2\sim$ the instantaneous bound. \mC\ the archived TFPT v4.5/v4.6 reported a clean
$A_s=2.124\times10^{-9}$ using the algebraic $N_\star=3/\phiz-1=55.42$; at that $N_\star$ the current normalisation
gives $2.047\times10^{-9}$ ($-1.9\sigma$) --- the old clean $A_s$ is near-instantaneous reheating, not the slow
channel. \mE\ nothing lost: the current theory replaced the posited algebraic $N_\star$ with one derived from
reheating physics and exposed honestly that $A_s$ requires fast (pre)heating. Record band $N_\star\in[50,60]$
(\vref{v84}) untouched.
\end{itemize}

\subsection*{2026-06-22 (the first-principles boundary, mapped: can $M_{\mathrm{Planck}}$ be computed? can $F_{\mathrm{transfer}}$ be first-principles? both NO, with the exact reason: \texttt{v339})}
\begin{itemize}
\item \textbf{\texttt{v339} (\texttt{FIRSTPRINCIPLES.BOUNDARY.01}).} An honest boundary audit answering two
head-on questions and fabricating nothing. \mE\ \emph{The Planck mass cannot be computed} ---
$v_{\mathrm{geo}}/M_{\mathrm{Planck}}$ is forbidden by theorem (No-Unit, \vref{v153}): a mass is mass-dimension
$1$ while every TFPT datum is dimension $0$, so no dimensionless map outputs an absolute scale. It is not a gap but
over-determined (gravity $=$ dark energy to $0.11\%$, \vref{v274}); every dimensionless \emph{ratio} is derived, only
the one unit is not. \mE\ \emph{$F_{\mathrm{transfer}}$ cannot be made first-principles} either: for $m_p/m_e$ the
structure is in place (\vref{v262}, carrier $b_3{=}{-}7$ $+$ closed electron) but two inputs stay external. \mX\
$\alpha_s(M_Z)$ is external \emph{because of a tension} --- the SM ($=$ TFPT content, no new states) does not unify
(1-loop spread ${\sim}8.8$ at $2{\times}10^{16}$ GeV, \vref{v246}). The lattice $C_p\!\sim\!2.83$ is parameter-free
but lattice-only (\vref{v35}); $\eta_B$/axion ride on cosmological initial conditions; Koide $Q{=}2/3$ is structural
\mE. \mE\ The residuals classify into four honest kinds: forbidden-by-theorem, external-and-tensioned,
external-but-parameter-free, cosmological. A classification, not a solution.
\end{itemize}

\subsection*{2026-06-22 (next-steps round: the one open lemma literature-anchored, the Decoupling Theorem, the holomorphy discriminator hardened in Lean, and the $\theta_{13}$ budget: \texttt{v336}, \texttt{v337}, \texttt{v338}, \texttt{FORM.CARTAN.DET.01})}
\begin{itemize}
\item \textbf{\texttt{v336} (\texttt{SEAM.EQUIV.CONTINUUM.01}) --- the one open lemma, literature-anchored.}
SEAM.EQUIV.01's only open input (the continuum OS reconstruction of the gapped quasi-free collar) is split into a
CITABLE part and a narrow open part. \mC\ continuum leg: Osborne--Stottmeister (\emph{CFT
from Lattice Fermions}, doi:10.1007/s00220-022-04521-8, CMP \textbf{398} (2023) 219) give, building on the
MMST (Morinelli--Morsella--Stottmeister--Tanimoto) operator-algebraic (wavelet) renormalization framework
(CMP 387 (2021) 299; PRL 127 (2021) 230601), the chiral-CFT scaling limit of free lattice fermions (Koo--Saleur
$\to$ Virasoro, correct $c$); WZW range $\mathrm{rank}\le c\le D$, and $(E_8)_1$ ($c{=}8$, rank 8, $D{=}16$
Majoranas, $SO(16)_1\!\to\!(E_8)_1$) is inside it. \mC\ OS leg: Adamo--Moriwaki--Tanimoto (arXiv:2407.18222, 2024)
prove the conformal OS axioms for unitary lattice VOAs. \mE\ target pinned by $c{=}\gcar{+}\Nfam{=}8$,
$\det\mathrm{Cartan}(E_8){=}1$. \mO\ narrowed residual $=$ ``the raw collar's massless scaling limit is the
holomorphic $(E_8)_1$ net'' ($\det K{=}1$), not a from-scratch construction.
\item \textbf{\texttt{v337} (\texttt{DECOUPLING.THEOREM.01}) --- the Decoupling Theorem.} \mE\ every dimensionless
readout factors through the gapped admissible transfer spectrum $\{1,(2/3)^6,(1/3)^6\}$; finite susceptibility
$\chi=1/(1-(2/3)^6)=729/665$ gap-suppresses the ambient contribution (margin $6\ln\tfrac32-31/(4\pi^2)\approx
1.648>0$, with $2\dim(E_8)c_3^2=31/(4\pi^2)$ exactly), so no readout depends on the ambient measure $\Rightarrow$
\texttt{QG.AMB.01} decoupled; negative control (gap closed $\Rightarrow\chi$ diverges) makes it non-vacuous. This is
the theorem that makes the \texttt{QG.AMB.01} reclassification (\vref{v335}) honest; consolidates
\vref{v76}/\vref{v311}/\vref{v330}/\vref{v332}.
\item \textbf{\texttt{v338} (\texttt{THETA13.BUDGET.01}) --- the $\theta_{13}$ budget, fenced.} The $+2\sigma$
pressure point (\vref{v328}) quantified against \emph{both} NuFIT 6.0 NO variants: $+2.0\sigma$ vs the no-SK
central ($0.02195$) but $+1.7\sigma$ vs the with-SK best fit ($0.02215\pm0.00057$), a ${\sim}0.3\sigma$ systematic
spread; \mE\ $0.023108$ lies inside the with-SK $3\sigma$ band (upper $0.02388$), an upper-edge pressure point not
a falsification; \mX\ JUNO's sub-1\% precision turns it into a ${\sim}6\sigma$ kill-or-confirm.
\item \textbf{Lean (\texttt{FORM.CARTAN.DET.01}).} \texttt{CartanDeterminants.lean} extended with the explicit
$D_8=SO(16)$ Cartan matrix and a fully-proven (kernel-only) holomorphy discriminator: $D_8$ is even and symmetric
just like $E_8$ (same $c{=}8$) but $4=\det\mathrm{Cartan}(D_8)$ is genuinely not a unit while $\det\mathrm{Cartan}
(E_8)$ is --- so unimodularity (one anyon) is the load-bearing selector. \texttt{lake build} green.
\end{itemize}

\subsection*{2026-06-22 (keystone unification + status reframe: ONE open theorem, ONE decoupled foreign problem: \texttt{v335}, \texttt{FORM.QGEO.BW.01} \texttt{qgeoSymIsCorollary})}
\begin{itemize}
\item \textbf{\texttt{v335} (\texttt{SEAM.EQUIV.UNIFY.01}) --- the keystone unification.} A bird's-eye
consolidation of the v308/v323/v325/v329/v333 arc that collapses the ``two open structural items'' narrative.
\mE\ \textbf{\texttt{QGEO.SYM.01} is a COROLLARY of \texttt{SEAM.EQUIV.01}} (no extra premise): a chiral conformal
net has, by axiom, a M\"obius-covariant vacuum (the unique invariant positive-energy vector); rotations $U(1)$ are
the compact subgroup of the M\"obius group, so the vacuum is rotation-invariant, and the $\mu_4$ clock $=$ rotation
by $\pi/2$ ($(e^{i\pi/2})^4{=}1$) fixes it $\Rightarrow\omega\circ\rho{=}\omega$. So \texttt{SEAM.EQUIV.01}
$\Rightarrow$ \texttt{QGEO.SYM.01} with NO extra premise (the v323 ``rotation-invariant vacuum'' is a net axiom)
--- the two bedrock items collapse to ONE. \mE\ \textbf{the one keystone theorem}: \texttt{SEAM.EQUIV.01} $=$ a
construction theorem ($c{=}\gcar{+}\Nfam{=}8$, $\det\mathrm{Cartan}(E_8){=}1$, gap $6\ln\tfrac32>0$ derived) with
EXACTLY one open continuum lemma (the rigorous continuum OS reconstruction $+$ invertibility of the raw collar).
\mE\ \textbf{\texttt{QG.AMB.01} decoupled, not a gate}: margin $\Delta-31/(4\pi^2)\approx1.648>0$
(v76/v311/v330/v332) $\Rightarrow$ no TFPT prediction depends on the ambient measure; it is the \emph{general}
Euclidean-QG conformal-factor problem (Gibbons--Hawking--Perry 1978), inherited not TFPT-specific --- reclassified
from ``TFPT structural frontier'' to ``decoupled general QG problem''. \mE\ \textbf{collapsed status}: open
structural items $2\to1$; live residual $=v_{\mathrm{geo}}+G_{\mathrm{net}}({=}\texttt{SEAM.EQUIV.01})+
F_{\mathrm{transfer}}$. \mO\ the one residual $=$ \texttt{SEAM.EQUIV.01}'s continuum OS reconstruction.
\item \textbf{Lean (\texttt{FORM.QGEO.BW.01}, \texttt{qgeoSymIsCorollary}).} \texttt{BWKeystone.lean} extended with the
cited conformal-net vacuum axiom \texttt{conformal\_net\_vacuum\_rotation\_invariant}; with it,
\texttt{StateInvariance} ($=$\texttt{QGEO.SYM.01}) follows from \texttt{SeamIsE8} alone, so \texttt{QGEO.SYM.01} is a
COROLLARY of \texttt{SEAM.EQUIV.01} with no independent open premise (\texttt{\#print axioms} confirms only the chain
$+$ cited BW/net axioms). \texttt{lake build} green (3362 jobs).
\item \textbf{Status deep-sync.} The ledger (\texttt{QG.AMB.01} reclassified to ``decoupled general QG; non-TFPT'';
\texttt{QGEO.SYM.01} marked a corollary; new \texttt{SEAM.EQUIV.UNIFY.01} row), the canonical status card
(\texttt{tfpt\_status.tex}), \texttt{introduction} (status core $+$ gate table), \texttt{origin\_theory},
\texttt{tfpt\_4\_frontier}, \texttt{tfpt\_5\_redteam}, \texttt{tfpt\_research\_contracts}, \texttt{README},
\texttt{next.txt} and the website (\texttt{OpenGates}, FAQ, \texttt{papers.ts}) all moved to ``one solvable theorem
(\texttt{SEAM.EQUIV.01}) $+$ one decoupled foreign problem (\texttt{QG.AMB.01})''. No falsifiable prediction
changed; this is a status-framing sharpening backed by \texttt{v335} $+$ the Lean corollary.
\end{itemize}

